\IfFileExists{stacks-project.cls}{%
\documentclass{stacks-project}
}{%
\documentclass{amsart}
}

% For dealing with references we use the comment environment
\usepackage{verbatim}
\newenvironment{reference}{\comment}{\endcomment}
%\newenvironment{reference}{}{}
\newenvironment{slogan}{\comment}{\endcomment}
\newenvironment{history}{\comment}{\endcomment}

% For commutative diagrams we use Xy-pic
\usepackage[all]{xy}

% We use 2cell for 2-commutative diagrams.
\xyoption{2cell}
\UseAllTwocells

% We use multicol for the list of chapters between chapters
\usepackage{multicol}

% This is generally recommended for better output
\usepackage{lmodern}
\usepackage[T1]{fontenc}

% For cross-file-references
\usepackage{xr-hyper}

% Package for hypertext links:
\usepackage{hyperref}

% For any local file, say "hello.tex" you want to link to please
% use \externaldocument[hello-]{hello}
\externaldocument[introduction-]{introduction}
\externaldocument[conventions-]{conventions}
\externaldocument[sets-]{sets}
\externaldocument[categories-]{categories}
\externaldocument[topology-]{topology}
\externaldocument[sheaves-]{sheaves}
\externaldocument[sites-]{sites}
\externaldocument[stacks-]{stacks}
\externaldocument[fields-]{fields}
\externaldocument[algebra-]{algebra}
\externaldocument[brauer-]{brauer}
\externaldocument[homology-]{homology}
\externaldocument[derived-]{derived}
\externaldocument[simplicial-]{simplicial}
\externaldocument[more-algebra-]{more-algebra}
\externaldocument[smoothing-]{smoothing}
\externaldocument[modules-]{modules}
\externaldocument[sites-modules-]{sites-modules}
\externaldocument[injectives-]{injectives}
\externaldocument[cohomology-]{cohomology}
\externaldocument[sites-cohomology-]{sites-cohomology}
\externaldocument[dga-]{dga}
\externaldocument[dpa-]{dpa}
\externaldocument[sdga-]{sdga}
\externaldocument[hypercovering-]{hypercovering}
\externaldocument[schemes-]{schemes}
\externaldocument[constructions-]{constructions}
\externaldocument[properties-]{properties}
\externaldocument[morphisms-]{morphisms}
\externaldocument[coherent-]{coherent}
\externaldocument[divisors-]{divisors}
\externaldocument[limits-]{limits}
\externaldocument[varieties-]{varieties}
\externaldocument[topologies-]{topologies}
\externaldocument[descent-]{descent}
\externaldocument[perfect-]{perfect}
\externaldocument[more-morphisms-]{more-morphisms}
\externaldocument[flat-]{flat}
\externaldocument[groupoids-]{groupoids}
\externaldocument[more-groupoids-]{more-groupoids}
\externaldocument[etale-]{etale}
\externaldocument[chow-]{chow}
\externaldocument[intersection-]{intersection}
\externaldocument[pic-]{pic}
\externaldocument[weil-]{weil}
\externaldocument[adequate-]{adequate}
\externaldocument[dualizing-]{dualizing}
\externaldocument[duality-]{duality}
\externaldocument[discriminant-]{discriminant}
\externaldocument[derham-]{derham}
\externaldocument[local-cohomology-]{local-cohomology}
\externaldocument[algebraization-]{algebraization}
\externaldocument[curves-]{curves}
\externaldocument[resolve-]{resolve}
\externaldocument[models-]{models}
\externaldocument[functors-]{functors}
\externaldocument[equiv-]{equiv}
\externaldocument[pione-]{pione}
\externaldocument[etale-cohomology-]{etale-cohomology}
\externaldocument[proetale-]{proetale}
\externaldocument[relative-cycles-]{relative-cycles}
\externaldocument[more-etale-]{more-etale}
\externaldocument[trace-]{trace}
\externaldocument[crystalline-]{crystalline}
\externaldocument[spaces-]{spaces}
\externaldocument[spaces-properties-]{spaces-properties}
\externaldocument[spaces-morphisms-]{spaces-morphisms}
\externaldocument[decent-spaces-]{decent-spaces}
\externaldocument[spaces-cohomology-]{spaces-cohomology}
\externaldocument[spaces-limits-]{spaces-limits}
\externaldocument[spaces-divisors-]{spaces-divisors}
\externaldocument[spaces-over-fields-]{spaces-over-fields}
\externaldocument[spaces-topologies-]{spaces-topologies}
\externaldocument[spaces-descent-]{spaces-descent}
\externaldocument[spaces-perfect-]{spaces-perfect}
\externaldocument[spaces-more-morphisms-]{spaces-more-morphisms}
\externaldocument[spaces-flat-]{spaces-flat}
\externaldocument[spaces-groupoids-]{spaces-groupoids}
\externaldocument[spaces-more-groupoids-]{spaces-more-groupoids}
\externaldocument[bootstrap-]{bootstrap}
\externaldocument[spaces-pushouts-]{spaces-pushouts}
\externaldocument[spaces-chow-]{spaces-chow}
\externaldocument[groupoids-quotients-]{groupoids-quotients}
\externaldocument[spaces-more-cohomology-]{spaces-more-cohomology}
\externaldocument[spaces-simplicial-]{spaces-simplicial}
\externaldocument[spaces-duality-]{spaces-duality}
\externaldocument[formal-spaces-]{formal-spaces}
\externaldocument[restricted-]{restricted}
\externaldocument[spaces-resolve-]{spaces-resolve}
\externaldocument[formal-defos-]{formal-defos}
\externaldocument[defos-]{defos}
\externaldocument[cotangent-]{cotangent}
\externaldocument[examples-defos-]{examples-defos}
\externaldocument[algebraic-]{algebraic}
\externaldocument[examples-stacks-]{examples-stacks}
\externaldocument[stacks-sheaves-]{stacks-sheaves}
\externaldocument[criteria-]{criteria}
\externaldocument[artin-]{artin}
\externaldocument[quot-]{quot}
\externaldocument[stacks-properties-]{stacks-properties}
\externaldocument[stacks-morphisms-]{stacks-morphisms}
\externaldocument[stacks-limits-]{stacks-limits}
\externaldocument[stacks-cohomology-]{stacks-cohomology}
\externaldocument[stacks-perfect-]{stacks-perfect}
\externaldocument[stacks-introduction-]{stacks-introduction}
\externaldocument[stacks-more-morphisms-]{stacks-more-morphisms}
\externaldocument[stacks-geometry-]{stacks-geometry}
\externaldocument[moduli-]{moduli}
\externaldocument[moduli-curves-]{moduli-curves}
\externaldocument[examples-]{examples}
\externaldocument[exercises-]{exercises}
\externaldocument[guide-]{guide}
\externaldocument[desirables-]{desirables}
\externaldocument[coding-]{coding}
\externaldocument[obsolete-]{obsolete}
\externaldocument[fdl-]{fdl}
\externaldocument[index-]{index}

% Theorem environments.
%
\theoremstyle{plain}
\newtheorem{theorem}[subsection]{Theorem}
\newtheorem{proposition}[subsection]{Proposition}
\newtheorem{lemma}[subsection]{Lemma}

\theoremstyle{definition}
\newtheorem{definition}[subsection]{Definition}
\newtheorem{example}[subsection]{Example}
\newtheorem{exercise}[subsection]{Exercise}
\newtheorem{situation}[subsection]{Situation}

\theoremstyle{remark}
\newtheorem{remark}[subsection]{Remark}
\newtheorem{remarks}[subsection]{Remarks}

\numberwithin{equation}{subsection}

% Macros
%
\def\lim{\mathop{\mathrm{lim}}\nolimits}
\def\colim{\mathop{\mathrm{colim}}\nolimits}
\def\Spec{\mathop{\mathrm{Spec}}}
\def\Hom{\mathop{\mathrm{Hom}}\nolimits}
\def\Ext{\mathop{\mathrm{Ext}}\nolimits}
\def\SheafHom{\mathop{\mathcal{H}\!\mathit{om}}\nolimits}
\def\SheafExt{\mathop{\mathcal{E}\!\mathit{xt}}\nolimits}
\def\Sch{\mathit{Sch}}
\def\Mor{\mathop{\mathrm{Mor}}\nolimits}
\def\Ob{\mathop{\mathrm{Ob}}\nolimits}
\def\Sh{\mathop{\mathit{Sh}}\nolimits}
\def\NL{\mathop{N\!L}\nolimits}
\def\CH{\mathop{\mathrm{CH}}\nolimits}
\def\proetale{{pro\text{-}\acute{e}tale}}
\def\etale{{\acute{e}tale}}
\def\QCoh{\mathit{QCoh}}
\def\Ker{\mathop{\mathrm{Ker}}}
\def\Im{\mathop{\mathrm{Im}}}
\def\Coker{\mathop{\mathrm{Coker}}}
\def\Coim{\mathop{\mathrm{Coim}}}

% Boxtimes
%
\DeclareMathSymbol{\boxtimes}{\mathbin}{AMSa}{"02}

%
% Macros for moduli stacks/spaces
%
\def\QCohstack{\mathcal{QC}\!\mathit{oh}}
\def\Cohstack{\mathcal{C}\!\mathit{oh}}
\def\Spacesstack{\mathcal{S}\!\mathit{paces}}
\def\Quotfunctor{\mathrm{Quot}}
\def\Hilbfunctor{\mathrm{Hilb}}
\def\Curvesstack{\mathcal{C}\!\mathit{urves}}
\def\Polarizedstack{\mathcal{P}\!\mathit{olarized}}
\def\Complexesstack{\mathcal{C}\!\mathit{omplexes}}
% \Pic is the operator that assigns to X its picard group, usage \Pic(X)
% \Picardstack_{X/B} denotes the Picard stack of X over B
% \Picardfunctor_{X/B} denotes the Picard functor of X over B
\def\Pic{\mathop{\mathrm{Pic}}\nolimits}
\def\Picardstack{\mathcal{P}\!\mathit{ic}}
\def\Picardfunctor{\mathrm{Pic}}
\def\Deformationcategory{\mathcal{D}\!\mathit{ef}}
\usepackage{cleveref}
\crefname{originallemma}{original lemma}{original lemmas}
\Crefname{originallemma}{Original lemma}{Original lemmas}
\makeatletter
\newlabel{sites-cohomology-lemma-pseudo-coherent-pullback}{{08H4}{source}{}{sourcepointer-1}{}}
\newlabel{sites-cohomology-lemma-pseudo-coherent-pullback@cref}{{[originallemma][0][]08H4}{[1][0][]source}{}{}{}}
\makeatother
\crefname{originallemma}{original lemma}{original lemmas}
\Crefname{originallemma}{Original lemma}{Original lemmas}
\makeatletter
\newlabel{spaces-perfect-lemma-pseudo-coherent}{{08JL}{source}{}{sourcepointer-2}{}}
\newlabel{spaces-perfect-lemma-pseudo-coherent@cref}{{[originallemma][0][]08JL}{[1][0][]source}{}{}{}}
\makeatother
\crefname{originalproposition}{original proposition}{original propositions}
\Crefname{originalproposition}{Original proposition}{Original propositions}
\makeatletter
\newlabel{spaces-descent-proposition-fpqc-descent-quasi-coherent}{{04W8}{source}{}{sourcepointer-3}{}}
\newlabel{spaces-descent-proposition-fpqc-descent-quasi-coherent@cref}{{[originalproposition][0][]04W8}{[1][0][]source}{}{}{}}
\makeatother
\crefname{originallemma}{original lemma}{original lemmas}
\Crefname{originallemma}{Original lemma}{Original lemmas}
\makeatletter
\newlabel{bootstrap-lemma-locally-algebraic-space-finite-type}{{04U0}{source}{}{sourcepointer-4}{}}
\newlabel{bootstrap-lemma-locally-algebraic-space-finite-type@cref}{{[originallemma][0][]04U0}{[1][0][]source}{}{}{}}
\makeatother
\crefname{originallemma}{original lemma}{original lemmas}
\Crefname{originallemma}{Original lemma}{Original lemmas}
\makeatletter
\newlabel{spaces-morphisms-lemma-affine-local}{{03WG}{source}{}{sourcepointer-5}{}}
\newlabel{spaces-morphisms-lemma-affine-local@cref}{{[originallemma][0][]03WG}{[1][0][]source}{}{}{}}
\makeatother
\crefname{originallemma}{original lemma}{original lemmas}
\Crefname{originallemma}{Original lemma}{Original lemmas}
\makeatletter
\newlabel{spaces-morphisms-lemma-finite-presentation-local}{{0410}{source}{}{sourcepointer-6}{}}
\newlabel{spaces-morphisms-lemma-finite-presentation-local@cref}{{[originallemma][0][]0410}{[1][0][]source}{}{}{}}
\makeatother
\crefname{originallemma}{original lemma}{original lemmas}
\Crefname{originallemma}{Original lemma}{Original lemmas}
\makeatletter
\newlabel{spaces-perfect-lemma-derived-quasi-coherent-small-etale-site}{{071Q}{source}{}{sourcepointer-7}{}}
\newlabel{spaces-perfect-lemma-derived-quasi-coherent-small-etale-site@cref}{{[originallemma][0][]071Q}{[1][0][]source}{}{}{}}
\makeatother
\crefname{originallemma}{original lemma}{original lemmas}
\Crefname{originallemma}{Original lemma}{Original lemmas}
\makeatletter
\newlabel{spaces-perfect-lemma-descend-pseudo-coherent}{{08HE}{source}{}{sourcepointer-8}{}}
\newlabel{spaces-perfect-lemma-descend-pseudo-coherent@cref}{{[originallemma][0][]08HE}{[1][0][]source}{}{}{}}
\makeatother
\crefname{originallemma}{original lemma}{original lemmas}
\Crefname{originallemma}{Original lemma}{Original lemmas}
\makeatletter
\newlabel{perfect-lemma-pseudo-coherent}{{08E5}{source}{}{sourcepointer-9}{}}
\newlabel{perfect-lemma-pseudo-coherent@cref}{{[originallemma][0][]08E5}{[1][0][]source}{}{}{}}
\makeatother
\crefname{originallemma}{original lemma}{original lemmas}
\Crefname{originallemma}{Original lemma}{Original lemmas}
\makeatletter
\newlabel{perfect-lemma-affine-compare-bounded}{{06Z0}{source}{}{sourcepointer-10}{}}
\newlabel{perfect-lemma-affine-compare-bounded@cref}{{[originallemma][0][]06Z0}{[1][0][]source}{}{}{}}
\makeatother
\crefname{originallemma}{original lemma}{original lemmas}
\Crefname{originallemma}{Original lemma}{Original lemmas}
\makeatletter
\newlabel{perfect-lemma-pseudo-coherent-affine}{{08E7}{source}{}{sourcepointer-11}{}}
\newlabel{perfect-lemma-pseudo-coherent-affine@cref}{{[originallemma][0][]08E7}{[1][0][]source}{}{}{}}
\makeatother
\crefname{originallemma}{original lemma}{original lemmas}
\Crefname{originallemma}{Original lemma}{Original lemmas}
\makeatletter
\newlabel{more-algebra-lemma-cut-complex-in-two}{{068U}{source}{}{sourcepointer-12}{}}
\newlabel{more-algebra-lemma-cut-complex-in-two@cref}{{[originallemma][0][]068U}{[1][0][]source}{}{}{}}
\makeatother
\crefname{originallemma}{original lemma}{original lemmas}
\Crefname{originallemma}{Original lemma}{Original lemmas}
\makeatletter
\newlabel{more-algebra-lemma-glue-perfect}{{066Y}{source}{}{sourcepointer-13}{}}
\newlabel{more-algebra-lemma-glue-perfect@cref}{{[originallemma][0][]066Y}{[1][0][]source}{}{}{}}
\makeatother
\crefname{originallemma}{original lemma}{original lemmas}
\Crefname{originallemma}{Original lemma}{Original lemmas}
\makeatletter
\newlabel{more-algebra-lemma-glue-tor-amplitude}{{066N}{source}{}{sourcepointer-14}{}}
\newlabel{more-algebra-lemma-glue-tor-amplitude@cref}{{[originallemma][0][]066N}{[1][0][]source}{}{}{}}
\makeatother
\crefname{originallemma}{original lemma}{original lemmas}
\Crefname{originallemma}{Original lemma}{Original lemmas}
\makeatletter
\newlabel{derived-lemma-negative-vanishing}{{05TC}{source}{}{sourcepointer-15}{}}
\newlabel{derived-lemma-negative-vanishing@cref}{{[originallemma][0][]05TC}{[1][0][]source}{}{}{}}
\makeatother
\crefname{originaldefinition}{original definition}{original definitions}
\Crefname{originaldefinition}{Original definition}{Original definitions}
\makeatletter
\newlabel{more-algebra-definition-pseudo-coherent}{{064Q}{source}{}{sourcepointer-16}{}}
\newlabel{more-algebra-definition-pseudo-coherent@cref}{{[originaldefinition][0][]064Q}{[1][0][]source}{}{}{}}
\makeatother
\crefname{originallemma}{original lemma}{original lemmas}
\Crefname{originallemma}{Original lemma}{Original lemmas}
\makeatletter
\newlabel{more-algebra-lemma-last-one-flat}{{0653}{source}{}{sourcepointer-17}{}}
\newlabel{more-algebra-lemma-last-one-flat@cref}{{[originallemma][0][]0653}{[1][0][]source}{}{}{}}
\makeatother
\crefname{originallemma}{original lemma}{original lemmas}
\Crefname{originallemma}{Original lemma}{Original lemmas}
\makeatletter
\newlabel{algebra-lemma-finite-projective}{{00NX}{source}{}{sourcepointer-18}{}}
\newlabel{algebra-lemma-finite-projective@cref}{{[originallemma][0][]00NX}{[1][0][]source}{}{}{}}
\makeatother
\crefname{originallemma}{original lemma}{original lemmas}
\Crefname{originallemma}{Original lemma}{Original lemmas}
\makeatletter
\newlabel{cohomology-lemma-bounded-flat-K-flat}{{06YD}{source}{}{sourcepointer-19}{}}
\newlabel{cohomology-lemma-bounded-flat-K-flat@cref}{{[originallemma][0][]06YD}{[1][0][]source}{}{}{}}
\makeatother
\title{Representability of the H\textasciicircum{}0 functor of a pseudo-coherent complex; perfectness under globally nonnegative fibres}
\author{The Stacks Project Authors}
\date{stacks-project:quot.tex:local-snapshot-20261012; GNU FDL 1.2 or later}
\begin{document}
\maketitle
\setcounter{subsection}{1}
\noindent\textit{Editorial note.} All delivered mathematical text is the author's own, copied from the pinned Stacks Project chapter file quot.tex; no statement or proof step was rewritten, reordered or removed. The counted span is the author environment and its own proof environment. References inside the span are handed to the assembler as navigation stubs that carry the canonical tag of the cited tag only; whether each of them is load-bearing is recorded on the stub itself.\par\medskip
\noindent\textit{Allowed standard prerequisites (editorial).}\begin{itemize}
\item Lemma sites-cohomology-lemma-pseudo-coherent-pullback (canonical tag 08H4). cited at proof lines 686-687. Lg\textasciicircum{}*K is pseudo-coherent on T.
\item Lemma spaces-perfect-lemma-pseudo-coherent (canonical tag 08JL). cited at proof lines 689-690 and 711-712. Cohomology sheaves of pseudo-coherent complexes are quasi-coherent; module translation.
\item Proposition spaces-descent-proposition-fpqc-descent-quasi-coherent (canonical tag 04W8). cited at proof lines 691-692. Cech description of sections of quasi-coherent sheaves.
\item Lemma bootstrap-lemma-locally-algebraic-space-finite-type (canonical tag 04U0). cited at proof line 701. Representability is etale local.
\item Lemma spaces-morphisms-lemma-affine-local (canonical tag 03WG). cited at proof line 705. Affineness is etale local.
\item Lemma spaces-morphisms-lemma-finite-presentation-local (canonical tag 0410). cited at proof lines 705-706. Finite presentation is etale local.
\item Lemma spaces-perfect-lemma-derived-quasi-coherent-small-etale-site (canonical tag 071Q). cited at proof line 713. Zariski/module translation of K.
\item Lemma spaces-perfect-lemma-descend-pseudo-coherent (canonical tag 08HE). cited at proof lines 713-714. Pseudo-coherence descends.
\item Lemma perfect-lemma-pseudo-coherent (canonical tag 08E5). cited at proof lines 719-721. M\textasciicircum{}bullet represents K by a complex of A-modules.
\item Lemma perfect-lemma-affine-compare-bounded (canonical tag 06Z0). cited at proof lines 720-721. Comparison of module and sheaf derived categories on affines.
\item Lemma perfect-lemma-pseudo-coherent-affine (canonical tag 08E7). cited at proof lines 721-722. Realization over the affine base.
\item Lemma more-algebra-lemma-cut-complex-in-two (canonical tag 068U). cited at proof lines 726-727. Canonical cut M\textasciicircum{}bullet = tau\_\{>=0\} + tau\_\{<=-1\}.
\item Lemma more-algebra-lemma-glue-perfect (canonical tag 066Y). cited at proof lines 733-734. Perfectness of the positive truncation.
\item Lemma more-algebra-lemma-glue-tor-amplitude (canonical tag 066N). cited at proof lines 734-735. Tor amplitude [0, b] of the positive truncation.
\item Lemma derived-lemma-negative-vanishing (canonical tag 05TC). cited at proof lines 742-743 (dual form). H\textasciicircum{}0(T, H\textasciicircum{}0(Lg\textasciicircum{}*K)) = H\textasciicircum{}0(T, Lg\textasciicircum{}*K).
\item Definition more-algebra-definition-pseudo-coherent (canonical tag 064Q). cited at proof lines 751-752. Bounded-above finite-free representative of pseudo-coherent complexes.
\item Lemma more-algebra-lemma-last-one-flat (canonical tag 0653). cited at proof lines 754-755. M = Coker(M\textasciicircum{}\{-1\} -> M\textasciicircum{}0) is flat.
\item Lemma algebra-lemma-finite-projective (canonical tag 00NX). cited at proof lines 755-756. Flat finite presentation gives finite locally free.
\item Lemma cohomology-lemma-bounded-flat-K-flat (canonical tag 06YD). cited at proof lines 761-762. The representative is K-flat, so derived pullback is term-wise.
\end{itemize}\medskip
\ifdefined\hypertarget\hypertarget{sourcepointer-1}{}\fi
\noindent\textit{Source reference.} Original lemma 08H4: sources/stacks-chow/tags.txt, source lines 10767--10767. Cited at proof lines 686-687: Lg\textasciicircum{}*K is pseudo-coherent on T; load-bearing.\par\medskip
\ifdefined\hypertarget\hypertarget{sourcepointer-2}{}\fi
\noindent\textit{Source reference.} Original lemma 08JL: sources/stacks-chow/tags.txt, source lines 10854--10854. Cited at proof lines 689-690 (cohomology sheaf quasi-coherent) and 711-712 (translation to modules); load-bearing twice.\par\medskip
\ifdefined\hypertarget\hypertarget{sourcepointer-3}{}\fi
\noindent\textit{Source reference.} Original proposition 04W8: sources/stacks-chow/tags.txt, source lines 6351--6351. Cited at proof lines 691-692: the Cech description of H\textasciicircum{}0(T, F); load-bearing for the sheaf property.\par\medskip
\ifdefined\hypertarget\hypertarget{sourcepointer-4}{}\fi
\noindent\textit{Source reference.} Original lemma 04U0: sources/stacks-chow/tags.txt, source lines 6273--6273. Cited at proof line 701: reduces representability to etale local on B; load-bearing.\par\medskip
\ifdefined\hypertarget\hypertarget{sourcepointer-5}{}\fi
\noindent\textit{Source reference.} Original lemma 03WG: sources/stacks-chow/tags.txt, source lines 5111--5111. Cited at proof line 705: affineness checked etale locally; load-bearing.\par\medskip
\ifdefined\hypertarget\hypertarget{sourcepointer-6}{}\fi
\noindent\textit{Source reference.} Original lemma 0410: sources/stacks-chow/tags.txt, source lines 5275--5275. Cited at proof lines 705-706: finite presentation checked etale locally; load-bearing.\par\medskip
\ifdefined\hypertarget\hypertarget{sourcepointer-7}{}\fi
\noindent\textit{Source reference.} Original lemma 071Q: sources/stacks-chow/tags.txt, source lines 8999--8999. Cited at proof line 713: lets the proof think of K on the Zariski topology; load-bearing.\par\medskip
\ifdefined\hypertarget\hypertarget{sourcepointer-8}{}\fi
\noindent\textit{Source reference.} Original lemma 08HE: sources/stacks-chow/tags.txt, source lines 10777--10777. Cited at proof lines 713-714: pseudo-coherence descends, completing the Zariski translation; load-bearing.\par\medskip
\ifdefined\hypertarget\hypertarget{sourcepointer-9}{}\fi
\noindent\textit{Source reference.} Original lemma 08E5: sources/stacks-chow/tags.txt, source lines 10661--10661. Cited at proof lines 719-721: existence of the pseudo-coherent complex of A-modules M\textasciicircum{}bullet; load-bearing.\par\medskip
\ifdefined\hypertarget\hypertarget{sourcepointer-10}{}\fi
\noindent\textit{Source reference.} Original lemma 06Z0: sources/stacks-chow/tags.txt, source lines 8904--8904. Cited at proof lines 720-721: K is the object corresponding to M\textasciicircum{}bullet; load-bearing.\par\medskip
\ifdefined\hypertarget\hypertarget{sourcepointer-11}{}\fi
\noindent\textit{Source reference.} Original lemma 08E7: sources/stacks-chow/tags.txt, source lines 10663--10663. Cited at proof lines 721-722: M\textasciicircum{}bullet realizes K over the affine base; load-bearing.\par\medskip
\ifdefined\hypertarget\hypertarget{sourcepointer-12}{}\fi
\noindent\textit{Source reference.} Original lemma 068U: sources/stacks-chow/tags.txt, source lines 8018--8018. Cited at proof lines 726-727: M\textasciicircum{}bullet = tau\_\{>=0\} + tau\_\{<=-1\} with the positive part perfect of tor amplitude [0, b]; load-bearing.\par\medskip
\ifdefined\hypertarget\hypertarget{sourcepointer-13}{}\fi
\noindent\textit{Source reference.} Original lemma 066Y: sources/stacks-chow/tags.txt, source lines 7952--7952. Cited at proof lines 733-734: perfectness of tau\_\{>=0\}M\textasciicircum{}bullet; load-bearing.\par\medskip
\ifdefined\hypertarget\hypertarget{sourcepointer-14}{}\fi
\noindent\textit{Source reference.} Original lemma 066N: sources/stacks-chow/tags.txt, source lines 7942--7942. Cited at proof lines 734-735: tor amplitude [0, b] of tau\_\{>=0\}M\textasciicircum{}bullet; load-bearing.\par\medskip
\ifdefined\hypertarget\hypertarget{sourcepointer-15}{}\fi
\noindent\textit{Source reference.} Original lemma 05TC: sources/stacks-chow/tags.txt, source lines 7476--7476. Cited at proof lines 742-743 (dual form): H\textasciicircum{}0(T, H\textasciicircum{}0(Lg\textasciicircum{}*K)) = H\textasciicircum{}0(T, Lg\textasciicircum{}*K); load-bearing for replacing K by tau\_\{>=0\}K.\par\medskip
\ifdefined\hypertarget\hypertarget{sourcepointer-16}{}\fi
\noindent\textit{Source reference.} Original definition 064Q: sources/stacks-chow/tags.txt, source lines 7874--7874. Cited at proof lines 751-752 (with following discussion): M\textasciicircum{}bullet may be taken bounded above of finite free modules; load-bearing.\par\medskip
\ifdefined\hypertarget\hypertarget{sourcepointer-17}{}\fi
\noindent\textit{Source reference.} Original lemma 0653: sources/stacks-chow/tags.txt, source lines 7887--7887. Cited at proof lines 754-755: M = Coker(M\textasciicircum{}\{-1\} -> M\textasciicircum{}0) is flat; load-bearing.\par\medskip
\ifdefined\hypertarget\hypertarget{sourcepointer-18}{}\fi
\noindent\textit{Source reference.} Original lemma 00NX: sources/stacks-chow/tags.txt, source lines 875--875. Cited at proof lines 755-756: flat finite M is finite locally free; load-bearing.\par\medskip
\ifdefined\hypertarget\hypertarget{sourcepointer-19}{}\fi
\noindent\textit{Source reference.} Original lemma 06YD: sources/stacks-chow/tags.txt, source lines 8882--8882. Cited at proof lines 761-762: the two-term-and-up complex is K-flat, so derived pullback is computed term-wise; load-bearing.\par\medskip
\begin{lemma}
\label{lemma-cohomology-perfect-complex}
Let $S$ be a scheme. Let $B$ be an algebraic space over $S$.
Let $K$ be a pseudo-coherent object of $D(\mathcal{O}_B)$.
\begin{enumerate}
\item If for all $g : T \to B$ in $(\Sch/B)$ the cohomology sheaf
$H^{-1}(Lg^*K)$ is zero, then the functor
$$
(\Sch/B)^{opp} \longrightarrow \textit{Sets},\quad
(g : T \to B) \longmapsto H^0(T, H^0(Lg^*K))
$$
is an algebraic space affine and of finite presentation over $B$.
\item If for all $g : T \to B$ in $(\Sch/B)$ the cohomology sheaves
$H^i(Lg^*K)$ are zero for $i < 0$, then $K$ is perfect,
$K$ locally has tor amplitude in $[0, b]$, and the functor
$$
(\Sch/B)^{opp} \longrightarrow \textit{Sets},\quad
(g : T \to B) \longmapsto H^0(T, Lg^*K)
$$
is an algebraic space affine and of finite presentation over $B$.
\end{enumerate}
\end{lemma}

\begin{proof}
Under the assumptions of (2) we have $H^0(T, Lg^*K) = H^0(T, H^0(Lg^*K))$.
Let us prove that the rule $T \mapsto H^0(T, H^0(Lg^*K))$ satisfies the
sheaf property for the fppf topology. To do this assume we have an
fppf covering $\{h_i : T_i \to T\}$ of a scheme $g : T \to B$ over $B$.
Set $g_i = g \circ h_i$. Note that since $h_i$ is flat, we have
$Lh_i^* = h_i^*$ and $h_i^*$ commutes with taking cohomology. Hence
$$
H^0(T_i, H^0(Lg_i^*K)) =
H^0(T_i, H^0(h_i^*Lg^*K)) =
H^0(T, h_i^*H^0(Lg^*K))
$$
Similarly for the pullback to $T_i \times_T T_j$.
Since $Lg^*K$ is a pseudo-coherent complex on $T$
(Cohomology on Sites, Lemma
\ref{sites-cohomology-lemma-pseudo-coherent-pullback})
the cohomology sheaf $\mathcal{F} = H^0(Lg^*K)$ is quasi-coherent
(Derived Categories of Spaces, Lemma
\ref{spaces-perfect-lemma-pseudo-coherent}).
Hence by Descent on Spaces, Proposition
\ref{spaces-descent-proposition-fpqc-descent-quasi-coherent}
we see that
$$
H^0(T, \mathcal{F}) = \Ker(
\prod H^0(T_i, h_i^*\mathcal{F}) \to
\prod H^0(T_i \times_T T_j, (T_i \times_T T_j \to T)^*\mathcal{F}))
$$
In this way we see that the rules in (1) and (2) satisfy
the sheaf property for fppf coverings. This means we may apply
Bootstrap, Lemma \ref{bootstrap-lemma-locally-algebraic-space-finite-type}
to see it suffices to prove the representability \'etale locally on $B$.
Moreover, we may check whether the end result is affine and
of finite presentation \'etale locally on $B$, see
Morphisms of Spaces, Lemmas \ref{spaces-morphisms-lemma-affine-local} and
\ref{spaces-morphisms-lemma-finite-presentation-local}.
Hence we may assume that $B$ is an affine scheme.

\medskip\noindent
Assume $B = \Spec(A)$ is an affine scheme. By the results of
Derived Categories of Spaces, Lemmas
\ref{spaces-perfect-lemma-pseudo-coherent},
\ref{spaces-perfect-lemma-derived-quasi-coherent-small-etale-site}, and
\ref{spaces-perfect-lemma-descend-pseudo-coherent}
we deduce that in the rest of the proof we may think of $K$ as a perfect
object of the derived category of complexes of modules on $B$
in the Zariski topology. By 
Derived Categories of Schemes, Lemmas
\ref{perfect-lemma-pseudo-coherent},
\ref{perfect-lemma-affine-compare-bounded}, and
\ref{perfect-lemma-pseudo-coherent-affine} we can find a pseudo-coherent
complex $M^\bullet$ of $A$-modules such that $K$ is the corresponding
object of $D(\mathcal{O}_B)$. Our assumption on pullbacks implies
that $M^\bullet \otimes^\mathbf{L}_A \kappa(\mathfrak p)$
has vanishing $H^{-1}$ for all primes $\mathfrak p \subset A$.
By More on Algebra, Lemma \ref{more-algebra-lemma-cut-complex-in-two}
we can write
$$
M^\bullet =
\tau_{\geq 0}M^\bullet \oplus \tau_{\leq - 1}M^\bullet
$$
with $\tau_{\geq 0}M^\bullet$ perfect with Tor amplitude in $[0, b]$
for some $b \geq 0$ (here we also have used
More on Algebra, Lemmas \ref{more-algebra-lemma-glue-perfect} and
\ref{more-algebra-lemma-glue-tor-amplitude}).
Note that in case (2) we also see that $\tau_{\leq - 1}M^\bullet = 0$
in $D(A)$ whence $M^\bullet$ and $K$ are perfect with
tor amplitude in $[0, b]$. For any $B$-scheme $g : T \to B$ we have
$$
H^0(T, H^0(Lg^*K)) = H^0(T, H^0(Lg^*\tau_{\geq 0}K))
$$
(by the dual of Derived Categories, Lemma
\ref{derived-lemma-negative-vanishing})
hence we may replace $K$ by $\tau_{\geq 0}K$ and correspondingly
$M^\bullet$ by $\tau_{\geq 0}M^\bullet$. In other words, we may
assume $M^\bullet$ has tor amplitude in $[0, b]$.

\medskip\noindent
Assume $M^\bullet$ has tor amplitude in $[0, b]$.
We may assume $M^\bullet$ is a bounded above complex of finite free
$A$-modules (by our definition of pseudo-coherent complexes, see
More on Algebra, Definition \ref{more-algebra-definition-pseudo-coherent}
and the discussion following the definition).
By More on Algebra, Lemma \ref{more-algebra-lemma-last-one-flat}
we see that $M = \Coker(M^{- 1} \to M^0)$ is flat. By
Algebra, Lemma \ref{algebra-lemma-finite-projective} we see that $M$
is finite locally free. Hence $M^\bullet$ is quasi-isomorphic to
$$
M \to M^1 \to M^2 \to \ldots \to M^d \to 0 \ldots
$$
Note that this is a K-flat complex
(Cohomology, Lemma \ref{cohomology-lemma-bounded-flat-K-flat}),
hence derived pullback of $K$ via a morphism $T \to B$ is computed
by the complex
$$
g^*\widetilde{M} \to g^*\widetilde{M^1} \to \ldots
$$
Thus it suffices to show that the functor
$$
(g : T \to B) \longmapsto
\Ker(
\Gamma(T,g^*\widetilde{M})
\to
\Gamma(T, g^*(\widetilde{M^1})
)
$$
is representable by an affine scheme of finite presentation over $B$.

\medskip\noindent
We may still replace $B$ by the members of an affine open covering
in order to prove this last statement. Hence we may assume that $M$
is finite free (recall that $M^1$ is finite free to begin with).
Write $M = A^{\oplus n}$ and $M^1 = A^{\oplus m}$. Let the map
$M \to M^1$ be given by the $m \times n$ matrix $(a_{ij})$ with
coefficients in $A$. Then $\widetilde{M} = \mathcal{O}_B^{\oplus n}$
and $\widetilde{M^1} = \mathcal{O}_B^{\oplus m}$. Thus the functor
above is equal to the functor
$$
(g : T \to B) \longmapsto
\{(f_1, \ldots, f_n) \in \Gamma(T, \mathcal{O}_T) \mid
\sum g^\sharp(a_{ij})f_i = 0,\ j = 1, \ldots, m\}
$$
Clearly this is representable by the affine scheme
$$
\Spec\left(A[x_1, \ldots, x_n]/(\sum a_{ij}x_i; j = 1, \ldots, m)\right)
$$
and the lemma has been proved.
\end{proof}

\end{document}
